\documentclass[proposal,indonesia]{ugmskripsi}
% Options:
%   skripsi / proposal        (default: skripsi)
%   satuspasi / duaspasi / satusetengahspasi (default: satusetengahspasi)
%   indonesia / inggris       (default: indonesia)

%------------------------------------------------------------
% Metadata Proposal
%------------------------------------------------------------
% TODO: Ganti metadata di bawah ini dengan data Anda sendiri
\titleind{Penerapan Pola Apps of Apps ArgoCD untuk Orkestrasi Deployment Layanan Repositori Ilmiah InvenioRDM pada Kubernetes}
\titleeng{Implementation of ArgoCD Apps of Apps Pattern for Orchestrating Scientific Repository Service InvenioRDM Deployment on Kubernetes}
\fullname{Muhammad Argya Vityasy}
\idnum{23/522547/PA/22475}
\degree{Sarjana}
\gelar{S.T.}
\yearsubmit{2026}
\program{Informatika}
\dept{Teknik Elektro dan Informatika}
\firstsupervisor{Guntur Budi Herwanto, S.Kom., M.Cs.}
\firstexaminer{Nama Dosen Penguji 1}
\secondexaminer{Nama Dosen Penguji 2}

%------------------------------------------------------------
% Packages
%------------------------------------------------------------
\usepackage[utf8]{inputenc}
\usepackage{graphicx}
\graphicspath{{figures/}}
\usepackage{amsmath}
\usepackage{amssymb}
\usepackage{booktabs}
\usepackage{longtable}
\usepackage{tabularx}
\usepackage{multirow}
\usepackage{float}
\usepackage{listings}
\usepackage{xcolor}
\usepackage{hyperref}
\usepackage{url}
\usepackage{enumitem}

% Code listing style
\lstset{
    basicstyle=\ttfamily\small,
    breaklines=true,
    frame=single,
    numbers=left,
    numberstyle=\tiny,
    keywordstyle=\color{blue},
    commentstyle=\color{gray},
    stringstyle=\color{orange},
    showstringspaces=false,
    tabsize=2
}

% Define YAML language for listings
\lstdefinelanguage{yaml}{
  keywords={true,false,null,y,n},
  sensitive=false,
  comment=[l]{\#},
  morecomment=[s]{/*}{*/},
  morestring=[b]',
  morestring=[b]"
}

% Define Docker language for listings
\lstdefinelanguage{docker}{
  alsoletter={-},
  keywords={FROM,RUN,CMD,LABEL,MAINTAINER,EXPOSE,ENV,ADD,COPY,ENTRYPOINT,VOLUME,USER,WORKDIR,ARG,ONBUILD,STOPSIGNAL,HEALTHCHECK,SHELL},
  sensitive=false,
  comment=[l]{\#},
  morestring=[b]',
  morestring=[b]"
}

%------------------------------------------------------------
% Document Body
%------------------------------------------------------------
\begin{document}

% Preliminary Pages (Proposal: minimal)
\cover
\titlepageind
\approvalpage

\tableofcontents

\abstractind
InvenioRDM merupakan platform repositori data penelitian open-source yang memiliki arsitektur microservices kompleks meliputi layanan web, worker, database, mesin pencari, dan penyimpanan objek. Deployment manual pada lingkungan Kubernetes memerlukan pengelolaan banyak komponen secara terpisah yang rentan terhadap kesalahan konfigurasi dan sulit untuk dilacak. Penelitian ini mengusulkan penerapan pola Apps of Apps pada ArgoCD untuk mengorkestrasi deployment seluruh komponen InvenioRDM secara deklaratif. Pendekatan GitOps yang digunakan memungkinkan seluruh state infrastruktur tersimpan dalam repositori Git sebagai single source of truth, sehingga deployment, pembaruan, dan rollback dapat dilakukan secara otomatis dan teraudit. Metode penelitian yang digunakan adalah Research and Development (R\&D) dengan tahapan analisis kebutuhan, perancangan arsitektur Apps of Apps, implementasi manifest ArgoCD, serta pengujian fungsional dan non-fungsional pada klaster Kubernetes. Hasil penelitian diharapkan dapat menghasilkan sistem deployment GitOps untuk InvenioRDM yang mampu mengurangi kompleksitas manajemen konfigurasi dan meningkatkan reliabilitas deployment.
\endabstractind

\abstracteng
InvenioRDM is an open-source research data repository platform with a complex microservices architecture including web services, workers, databases, search engines, and object storage. Manual deployment in Kubernetes environments requires managing many separate components, which is prone to configuration errors and difficult to track. This research proposes implementing the Apps of Apps pattern on ArgoCD to orchestrate the deployment of all InvenioRDM components declaratively. The GitOps approach enables the entire infrastructure state to be stored in a Git repository as a single source of truth, allowing automated and auditable deployments, updates, and rollbacks. The research method used is Research and Development (R\&D) with stages of requirements analysis, Apps of Apps architecture design, ArgoCD manifest implementation, and functional and non-functional testing on a Kubernetes cluster. The expected result is a GitOps deployment system for InvenioRDM that can reduce configuration management complexity and improve deployment reliability.
\endabstracteng

%------------------------------------------------------------
% Chapters
%------------------------------------------------------------
\chapter{PENDAHULUAN}

\section{Latar Belakang}

Dalam era digital, pengelolaan data penelitian menjadi aspek krusial bagi institusi akademik dan lembaga riset. Repositori data penelitian berfungsi sebagai wadah penyimpanan, publikasi, dan pengelolaan aset digital ilmiah yang memastikan keberlanjutan, aksesibilitas, dan citabilitas hasil penelitian. InvenioRDM merupakan salah satu platform repositori data penelitian open-source yang dikembangkan oleh CERN dan komunitas internasional. Platform ini menyediakan fitur lengkap mulai dari pengunggahan data, pemberian metadata, pencarian, hingga integrasi dengan sistem persistent identifier seperti DOI \citep{invenio2024}.

Arsitektur InvenioRDM bersifat microservices yang terdiri dari beberapa komponen independen yang saling berinteraksi. Komponen-komponen tersebut meliputi layanan web berbasis Flask, worker Celery untuk pemrosesan asinkron, database PostgreSQL untuk penyimpanan data, Elasticsearch untuk indeks pencarian, Redis untuk caching dan message broker, serta penyimpanan objek seperti MinIO atau Amazon S3. Kompleksitas arsitektur ini menuntut pendekatan deployment yang terstruktur dan dapat direproduksi dengan konsisten.

Kubernetes telah menjadi platform orkestrasi kontainer yang dominan dalam industri maupun akademisi. Kubernetes menyediakan abstraksi untuk deployment, scaling, dan manajemen aplikasi kontainerisasi. Namun, deployment aplikasi kompleks seperti InvenioRDM pada Kubernetes menghadapi tantangan tersendiri. Pengelolaan banyak komponen secara manual melalui perintah \texttt{kubectl} atau skrip \textit{ad-hoc} rentan terhadap \textit{configuration drift}, sulit dilacak, dan tidak memiliki mekanisme rollback yang terstandarisasi \citep{kubernetes2024}.

GitOps merupakan paradigma deployment modern yang mengusulkan penggunaan repositori Git sebagai \textit{single source of truth} untuk seluruh state infrastruktur dan aplikasi. Dalam GitOps, setiap perubahan pada sistem dilakukan melalui commit ke repositori Git, yang kemudian secara otomatis disinkronkan ke lingkungan target oleh agen GitOps. Pendekatan ini menawarkan beberapa keuntungan: deklaratif, terverifikasi, teraudit, dan mampu melakukan \textit{self-healing} \citep{gitops2023}.

ArgoCD adalah salah satu implementasi GitOps yang paling banyak digunakan untuk ekosistem Kubernetes. ArgoCD memungkinkan definisi aplikasi Kubernetes dalam bentuk manifest YAML atau Helm chart yang disimpan di Git, kemudian secara kontinu memantau dan menyelaraskan state klaster dengan state yang didefinisikan di Git. Salah satu pola yang didukung ArgoCD adalah \textit{Apps of Apps}, yaitu sebuah aplikasi ArgoCD induk yang mengelola beberapa aplikasi anak. Pola ini sangat cocok untuk aplikasi kompleks dengan banyak komponen seperti InvenioRDM \citep{argocd2024}.

Berdasarkan pemaparan di atas, terdapat kesenjangan antara kompleksitas deployment InvenioRDM dan ketersediaan solusi GitOps yang terintegrasi untuk platform tersebut. Penelitian ini bertujuan untuk mengisi kesenjangan tersebut dengan merancang dan mengimplementasikan sistem deployment berbasis GitOps menggunakan pola Apps of Apps ArgoCD untuk InvenioRDM pada klaster Kubernetes.

\section{Rumusan Masalah}

Berdasarkan latar belakang di atas, rumusan masalah dalam penelitian ini adalah:

\begin{enumerate}
  \item Bagaimana merancang arsitektur deployment GitOps menggunakan pola Apps of Apps ArgoCD untuk orkestrasi komponen-komponen InvenioRDM pada klaster Kubernetes?
  \item Bagaimana mengimplementasikan sistem deployment yang mampu mengelola dependensi antar komponen InvenioRDM secara deklaratif dan otomatis?
  \item Bagaimana mengukur efektivitas sistem deployment GitOps yang dikembangkan dibandingkan dengan pendekatan deployment manual pada Kubernetes?
\end{enumerate}

\section{Batasan Masalah}

Agar penelitian ini tetap fokus dan terukur, batasan masalah yang ditetapkan adalah:

\begin{enumerate}
  \item Penelitian dilakukan pada lingkungan klaster Kubernetes tunggal (single cluster) tanpa mempertimbangkan skenario multi-cluster.
  \item Implementasi GitOps menggunakan ArgoCD sebagai tools utama, tidak membandingkan dengan tools GitOps lain seperti Flux atau Rancher Fleet.
  \item Versi InvenioRDM yang digunakan adalah versi 12.x dengan konfigurasi standar yang mendukung komponen web, worker, database PostgreSQL, Elasticsearch, Redis, dan penyimpanan objek MinIO.
  \item Aspek migrasi data dari sistem lama ke InvenioRDM tidak termasuk dalam ruang lingkup penelitian.
  \item Pengujian difokuskan pada aspek fungsionalitas deployment, sinkronisasi otomatis, drift detection, dan rollback, tanpa melakukan pengujian beban (load testing) skala besar.
\end{enumerate}

\section{Tujuan Penelitian}

Tujuan dari penelitian ini adalah:

\begin{enumerate}
  \item Merancang arsitektur deployment GitOps menggunakan pola Apps of Apps ArgoCD untuk mengorkestrasi seluruh komponen InvenioRDM pada klaster Kubernetes.
  \item Mengimplementasikan sistem deployment yang mampu mengelola dependensi antar komponen InvenioRDM secara deklaratif melalui manifest ArgoCD dan Helm chart.
  \item Mengevaluasi efektivitas sistem deployment GitOps yang dikembangkan berdasarkan metrik deployment time, sinkronisasi otomatis, dan kemampuan rollback.
\end{enumerate}

\section{Manfaat Penelitian}

Penelitian ini diharapkan memberikan manfaat berikut:

\begin{enumerate}
  \item \textbf{Teoritis:} Menambah khazanah keilmuan di bidang DevOps dan GitOps, khususnya penerapan pola Apps of Apps untuk aplikasi microservices kompleks pada Kubernetes.
  \item \textbf{Praktis:} Menghasilkan sistem deployment terstandarisasi untuk InvenioRDM yang dapat digunakan oleh institusi akademik dalam mengelola repositori data penelitian dengan lebih andal, teraudit, dan mudah direproduksi.
\end{enumerate}

\section{Sistematika Penulisan}

Sistematika penulisan proposal tugas akhir ini adalah sebagai berikut:

\begin{itemize}
  \item \textbf{Bab I:} Pendahuluan --- berisi latar belakang, rumusan masalah, batasan masalah, tujuan, manfaat, dan sistematika penulisan.
  \item \textbf{Bab II:} Penelitian Terkait --- membahas state of the art mengenai InvenioRDM, Kubernetes, GitOps, ArgoCD, dan pola Apps of Apps, serta hasil-hasil penelitian sebelumnya yang relevan.
  \item \textbf{Bab III:} Metode dan Rancangan Penelitian --- menjelaskan metodologi Research and Development (R\&D), perancangan arsitektur sistem, bagan alir metodologi, lingkungan pengembangan, dan metode pengujian.
  \item \textbf{Bab IV:} Jadwal Penelitian --- memuat rencana jadwal pelaksanaan penelitian berbasis bulan dalam bentuk Gantt Chart.
\end{itemize}

\chapter{PENELITIAN TERKAIT}

\section{InvenioRDM sebagai Platform Repositori Data Penelitian}

InvenioRDM adalah platform repositori data penelitian yang dibangun di atas framework Invenio yang dikembangkan oleh CERN sejak tahun 2016. Platform ini dirancang untuk memenuhi kebutuhan institusi akademik dalam mengelola, mempublikasikan, dan memelihara data penelitian sesuai dengan prinsip FAIR (Findable, Accessible, Interoperable, Reusable) \citep{invenio2024}. InvenioRDM menyediakan fitur lengkap termasuk pengunggahan data dengan metadata terstruktur, pencarian full-text berbasis Elasticsearch, versioning dataset, kontrol akses berbasis peran, serta integrasi dengan sistem DOI melalui DataCite.

Arsitektur InvenioRDM mengikuti pola microservices yang memisahkan tanggung jawab ke dalam beberapa layanan independen \citep{invenio2024}. Komponen utama InvenioRDM meliputi: (1) \textbf{Web Application}, yaitu aplikasi Flask yang menangani HTTP request dan menyajikan antarmuka pengguna; (2) \textbf{Celery Workers}, yaitu worker yang menangani tugas-tugas asinkron seperti pemrosesan file, pengiriman email, dan indexing; (3) \textbf{PostgreSQL}, yaitu database relasional untuk menyimpan metadata record, user, dan konfigurasi; (4) \textbf{Elasticsearch}, yaitu mesin pencarian dan analitik untuk indexing dan query data; (5) \textbf{Redis}, yang berfungsi sebagai message broker untuk Celery dan caching layer; serta (6) \textbf{Object Storage}, seperti MinIO atau Amazon S3 untuk menyimpan file upload pengguna. Kompleksitas interaksi antar komponen ini menuntut pendekatan deployment yang terstruktur.

\section{Kubernetes sebagai Platform Orkestrasi}

Kubernetes adalah sistem orkestrasi kontainer open-source yang dikembangkan oleh Google dan saat ini dikelola oleh Cloud Native Computing Foundation (CNCF). Kubernetes menyediakan platform untuk automasi deployment, scaling, dan operasi aplikasi kontainerisasi \citep{kubernetes2024}. Abstraksi utama dalam Kubernetes meliputi Pod sebagai unit deploymen terkecil, Deployment untuk manajemen replika aplikasi stateless, StatefulSet untuk aplikasi stateful, Service untuk eksposur jaringan, ConfigMap dan Secret untuk manajemen konfigurasi, serta Ingress untuk routing HTTP/HTTPS.

Dalam konteks deployment aplikasi ilmiah dan akademik, Kubernetes telah banyak diadopsi untuk menjalankan platform data-intensive. Kemampuan Kubernetes dalam mengelola resource, melakukan self-healing, dan mendukung horizontal scaling menjadikannya pilihan utama untuk deployment platform seperti InvenioRDM yang memerlukan availability tinggi \citep{kubernetes2024}. Namun, pengelolaan manifest Kubernetes yang banyak dan kompleks untuk aplikasi multi-komponen seringkali menjadi tantangan tersendiri.

\section{GitOps sebagai Paradigma Deployment Modern}

GitOps adalah paradigma untuk mengelola infrastruktur dan aplikasi yang menggunakan repositori Git sebagai sumber kebenaran tunggal (single source of truth) \citep{gitops2023}. Konsep GitOps pertama kali dipopulerkan oleh Weaveworks pada tahun 2017 dan sejak itu menjadi standar de facto untuk deployment pada lingkungan cloud-native. Prinsip utama GitOps mencakup: (1) seluruh sistem dideklarasikan secara deklaratif; (2) state sistem yang diinginkan tersimpan dalam versi di Git; (3) perubahan state diterapkan secara otomatis melalui agen software; dan (4) agen software secara kontinu mendeteksi dan mengoreksi deviasi (drift) dari state yang diinginkan.

Pendekatan GitOps menawarkan beberapa keunggulan signifikan dibandingkan deployment imperatif tradisional. Pertama, setiap perubahan pada infrastruktur tercatat dalam history Git, memungkinkan audit trail yang lengkap. Kedua, mekanisme pull request dan code review dapat diterapkan pada perubahan infrastruktur, meningkatkan kualitas dan keamanan. Ketiga, rollback ke state sebelumnya dapat dilakukan dengan mudah melalui Git revert. Keempat, konfigurasi dapat direproduksi dengan konsisten di berbagai lingkungan \citep{gitops2023}.

\section{ArgoCD dan Pola Apps of Apps}

ArgoCD adalah tools GitOps yang dirancang khusus untuk ekosistem Kubernetes. ArgoCD mengimplementasikan model pull-based GitOps di mana agen ArgoCD yang berjalan di dalam klaster Kubernetes secara periodik memeriksa repositori Git dan menyelaraskan state klaster dengan state yang didefinisikan di Git \citep{argocd2024}. ArgoCD mendukung berbagai format manifest termasuk plain YAML, Kustomize, Helm, dan Jsonnet.

\textit{Apps of Apps} adalah salah satu pola arsitektur yang didukung ArgoCD untuk mengelola banyak aplikasi secara hierarkis. Dalam pola ini, sebuah aplikasi ArgoCD induk (parent app) didefinisikan untuk mengelola beberapa aplikasi anak (child apps). Aplikasi induk berisi manifest yang mereferensikan aplikasi-anak, sehingga dengan mendeploy aplikasi induk, seluruh aplikasi anak akan secara otomatis dibuat dan dikelola \citep{argocd2024}. Pola ini sangat bermanfaat untuk aplikasi yang terdiri dari banyak komponen dengan dependensi antar komponen, seperti halnya InvenioRDM.

Keunggulan pola Apps of Apps meliputi: (1) \textbf{Organisasi hierarkis}, memungkinkan pengelompokan aplikasi berdasarkan domain atau tim; (2) \textbf{Dependency management}, aplikasi anak dapat dikonfigurasi untuk menunggu aplikasi lain siap melalui mekanisme sync wave; (3) \textbf{Single entry point}, operator cukup mengelola satu aplikasi induk untuk mengontrol seluruh stack; dan (4) \textbf{Reusability}, struktur aplikasi induk dapat digunakan sebagai template untuk lingkungan berbeda (development, staging, production).

\section{Penelitian Terdahulu}

Beberapa penelitian telah dilakukan terkait penerapan GitOps pada berbagai konteks. \citep{limoncelli2018} membahas praktik terbaik untuk manajemen infrastruktur cloud yang menekankan pentingnya otomasi dan version control. \citep{gitops2023} menganalisis evolusi GitOps dari konsep hingga adopsi industri, termasuk perbandingan antara model pull-based dan push-based. Dalam konteks Kubernetes, \citep{kubernetes2024} memberikan panduan komprehensif untuk deployment aplikasi kompleks menggunakan Helm dan Kustomize.

Penelitian spesifik mengenai deployment platform repository pernah dilakukan oleh \citep{zhang2022} yang mengimplementasikan arsitektur microservices untuk platform data repository menggunakan Docker Compose. Namun, penelitian tersebut tidak memanfaatkan Kubernetes dan pendekatan GitOps. \citep{wang2021} membahas penggunaan ArgoCD untuk deployment aplikasi enterprise, namun tidak secara spesifik membahas pola Apps of Apps untuk aplikasi dengan banyak komponen dependen.

Berdasarkan tinjauan literatur di atas, belum ditemukan penelitian yang secara khusus merancang dan mengimplementasikan pola Apps of Apps ArgoCD untuk deployment InvenioRDM pada Kubernetes. Penelitian ini diharapkan dapat mengisi gap tersebut dengan memberikan solusi praktis yang terintegrasi.

\section{State of the Art}

Berdasarkan tinjauan pustaka yang telah dilakukan, state of the art dalam penelitian ini dapat dirangkum sebagai berikut:

\begin{enumerate}
  \item InvenioRDM memiliki arsitektur microservices yang kompleks dengan banyak komponen dependen, sehingga memerlukan pendekatan deployment terstruktur.
  \item Kubernetes menyediakan platform orkestrasi yang kuat, namun pengelolaan manifest untuk aplikasi multi-komponen masih menjadi tantangan.
  \item GitOps, khususnya melalui ArgoCD, menawarkan paradigma deployment deklaratif yang mampu mengatasi tantangan pengelolaan konfigurasi pada Kubernetes.
  \item Pola Apps of Apps merupakan solusi yang cocok untuk mengelola deployment aplikasi dengan banyak komponen dependen seperti InvenioRDM, namun implementasinya untuk platform ini belum pernah dilakukan.
\end{enumerate}

\chapter{METODE DAN RANCANGAN PENELITIAN}

\section{Deskripsi Umum Penelitian}

Penelitian ini mengusulkan pengembangan sistem deployment berbasis GitOps untuk platform repositori data penelitian InvenioRDM pada lingkungan klaster Kubernetes. Sistem yang dikembangkan menggunakan pola Apps of Apps pada ArgoCD untuk mengorkestrasi deployment seluruh komponen InvenioRDM secara deklaratif. Pendekatan Research and Development (R\&D) dipilih karena penelitian berfokus pada pengembangan produk berupa sistem deployment yang dapat diimplementasikan dan dievaluasi.

Secara umum, penelitian ini terdiri dari empat tahapan utama: (1) analisis kebutuhan komponen InvenioRDM dan karakteristik deployment pada Kubernetes; (2) perancangan arsitektur GitOps menggunakan pola Apps of Apps ArgoCD; (3) implementasi manifest ArgoCD, Helm chart, dan konfigurasi Kubernetes; serta (4) pengujian fungsionalitas dan evaluasi kinerja sistem deployment.

\section{Metodologi}

Metode penelitian yang digunakan adalah \textbf{Research and Development (R\&D)}. Metode R\&D dipilih karena penelitian ini bertujuan untuk menghasilkan produk berupa sistem deployment GitOps yang baru dan dapat diuji keefektifannya \citep{limoncelli2018}. Tahapan penelitian R\&D dalam konteks ini mengikuti siklus pengembangan sistem yang meliputi analisis, desain, implementasi, dan evaluasi.

Tahapan penelitian dilakukan secara sistematis sebagai berikut:

\begin{enumerate}
    \item \textbf{Tahap 1: Studi Literatur dan Analisis Kebutuhan}
    
    Tahap ini meliputi pengumpulan informasi mengenai arsitektur InvenioRDM, komponen-komponen yang diperlukan, dependensi antar komponen, serta kebutuhan resource pada Kubernetes. Studi literatur juga mencakup analisis dokumentasi ArgoCD, pola Apps of Apps, dan best practices GitOps deployment.
    
    \item \textbf{Tahap 2: Perancangan Arsitektur Sistem}
    
    Berdasarkan hasil analisis kebutuhan, dirancang arsitektur sistem deployment GitOps yang menggunakan pola Apps of Apps. Perancangan mencakup struktur repositori Git, hierarki aplikasi ArgoCD, konfigurasi Helm chart untuk masing-masing komponen, serta definisi resource Kubernetes (Deployment, Service, Ingress, PVC, Secret, ConfigMap).
    
    \item \textbf{Tahap 3: Implementasi Sistem}
    
    Tahap implementasi meliputi pembuatan klaster Kubernetes (menggunakan Minikube atau k3s), instalasi ArgoCD, pembuatan repositori Git untuk menyimpan manifest, pembuatan aplikasi ArgoCD induk dan anak, konfigurasi Helm chart untuk InvenioRDM beserta dependensinya, serta konfigurasi networking dan penyimpanan pada Kubernetes.
    
    \item \textbf{Tahap 4: Pengujian dan Evaluasi}
    
    Sistem yang telah diimplementasikan diuji menggunakan skenario pengujian yang telah dirancang. Metrik evaluasi meliputi waktu deployment, keberhasilan sinkronisasi otomatis, deteksi dan koreksi drift, serta kemampuan rollback. Hasil pengujian dianalisis untuk menentukan efektivitas sistem.
    
    \item \textbf{Tahap 5: Dokumentasi dan Pelaporan}
    
    Tahap akhir meliputi dokumentasi sistem, penulisan laporan penelitian, serta publikasi hasil jika memungkinkan.
\end{enumerate}

\section{Perancangan Sistem}

\subsection{Arsitektur Umum}

Arsitektur sistem yang dirancang terdiri dari empat lapisan utama yang saling berinteraksi:

\begin{itemize}
    \item \textbf{Lapisan 1: Git Repository} --- Berfungsi sebagai single source of truth yang menyimpan seluruh manifest Kubernetes, Helm values, dan konfigurasi ArgoCD dalam bentuk Infrastructure as Code (IaC).
    
    \item \textbf{Lapisan 2: ArgoCD} --- Agen GitOps yang berjalan di dalam klaster Kubernetes. ArgoCD secara kontinu memantau repositori Git dan menyelaraskan state klaster dengan state yang didefinisikan di Git. ArgoCD juga menyediakan antarmuka web untuk visualisasi dan manajemen aplikasi.
    
    \item \textbf{Lapisan 3: Kubernetes Cluster} --- Platform orkestrasi yang menjalankan seluruh komponen InvenioRDM dalam bentuk Pod, Deployment, Service, dan resource Kubernetes lainnya.
    
    \item \textbf{Lapisan 4: InvenioRDM Services} --- Kumpulan layanan aplikasi yang terdiri dari web application, Celery workers, PostgreSQL, Elasticsearch, Redis, dan MinIO object storage.
\end{itemize}

Dalam pola Apps of Apps, sebuah aplikasi ArgoCD induk bernama \texttt{inveniordm-platform} didefinisikan untuk mengelola seluruh stack. Aplikasi induk ini mereferensikan beberapa aplikasi anak yang masing-masing bertanggung jawab atas satu komponen utama:

\begin{itemize}
    \item \texttt{inveniordm-infra} --- Mengelola komponen infrastruktur (PostgreSQL, Elasticsearch, Redis, MinIO).
    \item \texttt{inveniordm-web} --- Mengelola deployment aplikasi web InvenioRDM.
    \item \texttt{inveniordm-worker} --- Mengelola deployment Celery workers.
\end{itemize}

Pembagian ini memungkinkan pengelolaan yang modular dan memfasilitasi penggunaan \textit{sync waves} untuk mengatur urutan deployment berdasarkan dependensi antar komponen.

% TODO: Ganti dengan diagram arsitektur sistem Anda.
% \begin{figure}[H]
%   \centering
%   \includegraphics[width=0.9\textwidth]{figures/arsitektur-sistem}
%   \caption{Diagram Arsitektur Sistem GitOps untuk InvenioRDM}
%   \label{fig:arsitektur}
% \end{figure}

\subsection{Bagan Alir Metodologi}

% TODO: Ganti dengan bagan alir metodologi Anda.
% \begin{figure}[H]
%   \centering
%   \includegraphics[width=0.8\textwidth]{figures/bagan-alir}
%   \caption{Bagan Alir Metodologi Penelitian}
%   \label{fig:bagan-alir}
% \end{figure}

Bagan alir metodologi penelitian menggambarkan alur dari tahap studi literatur hingga evaluasi. Setiap tahap memiliki deliverable yang jelas dan kriteria kelulusan sebelum melanjutkan ke tahap berikutnya.

\section{Lingkungan Pengembangan}

Lingkungan pengembangan yang digunakan dalam penelitian ini adalah:

\begin{itemize}
    \item \textbf{Sistem Operasi:} Ubuntu 22.04 LTS (untuk klaster Kubernetes) dan macOS/Linux/Windows dengan WSL2 (untuk development).
    \item \textbf{Container Runtime:} Docker Engine v24.x atau containerd.
    \item \textbf{Kubernetes:} Minikube v1.33+ atau k3s v1.29+ untuk lingkungan lokal.
    \item \textbf{GitOps Tools:} ArgoCD v2.10+.
    \item \textbf{Package Manager:} Helm v3.14+ untuk templating manifest Kubernetes.
    \item \textbf{Version Control:} Git dengan repositori pada GitHub.
    \item \textbf{Perangkat Keras:} Workstation dengan minimum 16 GB RAM, 8 core CPU, dan 100 GB storage (untuk menjalankan klaster Kubernetes lokal dengan banyak komponen).
\end{itemize}

\section{Metode Pengujian}

Pengujian dilakukan dengan beberapa pendekatan untuk memastikan sistem berfungsi sesuai harapan:

\begin{enumerate}
    \item \textbf{Pengujian Fungsional:}
    \begin{itemize}
        \item Deployment awal (fresh install) dari repositori Git ke klaster Kubernetes kosong.
        \item Verifikasi seluruh komponen InvenioRDM berjalan dan dapat berkomunikasi.
        \item Pengujian akses antarmuka web InvenioRDM melalui Ingress.
    \end{itemize}
    
    \item \textbf{Pengujian GitOps:}
    \begin{itemize}
        \item Perubahan konfigurasi pada Git dan verifikasi sinkronisasi otomatis oleh ArgoCD.
        \item Simulasi drift (perubahan manual pada klaster) dan verifikasi self-healing oleh ArgoCD.
        \item Pengujian rollback ke versi sebelumnya melalui Git revert.
    \end{itemize}
    
    \item \textbf{Pengujian Non-Fungsional:}
    \begin{itemize}
        \item Pengukuran waktu deployment dari awal hingga aplikasi siap digunakan.
        \item Pengukuran waktu sinkronisasi setelah perubahan pada Git.
        \item Evaluasi konsumsi resource (CPU, memory) pada klaster Kubernetes.
    \end{itemize}
\end{enumerate}

\chapter{JADWAL PENELITIAN}

Penelitian ini direncanakan untuk dilaksanakan selama enam bulan, dari bulan Februari 2026 hingga Juli 2026. Jadwal penelitian disusun berdasarkan tahapan metodologi Research and Development (R\&D) yang telah dijelaskan pada Bab III.

Tabel \ref{tab:jadwal} menunjukkan rencana jadwal pelaksanaan penelitian berbasis bulan. Simbol $\bullet$ menunjukkan bahwa kegiatan dilaksanakan pada bulan tersebut.

\begin{table}[H]
  \centering
  \caption{Jadwal Penelitian}
  \label{tab:jadwal}
  \begin{tabular}{|l|c|c|c|c|c|c|}
    \hline
    \textbf{Kegiatan} & \textbf{Feb} & \textbf{Mar} & \textbf{Apr} & \textbf{Mei} & \textbf{Jun} & \textbf{Jul} \\
    \hline
    Studi Literatur & $\bullet$ & $\bullet$ & & & & \\
    \hline
    Analisis Kebutuhan & $\bullet$ & $\bullet$ & & & & \\
    \hline
    Perancangan Sistem & & $\bullet$ & $\bullet$ & & & \\
    \hline
    Implementasi Sistem & & & $\bullet$ & $\bullet$ & & \\
    \hline
    Pengujian dan Evaluasi & & & & $\bullet$ & $\bullet$ & \\
    \hline
    Dokumentasi dan Penulisan & & $\bullet$ & $\bullet$ & $\bullet$ & $\bullet$ & $\bullet$ \\
    \hline
    Seminar Proposal / Sidang & & & & & & $\bullet$ \\
    \hline
  \end{tabular}
\end{table}

\section{Penjelasan Kegiatan}

\begin{enumerate}
    \item \textbf{Studi Literatur (Februari -- Maret 2026)}
    
    Kegiatan ini mencakup pengumpulan dan analisis literatur mengenai InvenioRDM, Kubernetes, GitOps, ArgoCD, dan pola Apps of Apps. Sumber literatur meliputi dokumentasi resmi, buku teks, jurnal ilmiah, dan artikel teknis.
    
    \item \textbf{Analisis Kebutuhan (Februari -- Maret 2026)}
    
    Tahap ini meliputi analisis kebutuhan komponen InvenioRDM, dependensi antar komponen, kebutuhan resource Kubernetes, serta identifikasi tools dan teknologi yang diperlukan untuk implementasi.
    
    \item \textbf{Perancangan Sistem (Maret -- April 2026)}
    
    Kegiatan perancangan meliputi pembuatan diagram arsitektur sistem, desain struktur repositori Git, perancangan hierarki Apps of Apps ArgoCD, dan perancangan manifest Kubernetes serta Helm chart.
    
    \item \textbf{Implementasi Sistem (April -- Mei 2026)}
    
    Tahap implementasi meliputi pembuatan klaster Kubernetes, instalasi ArgoCD, pembuatan repositori Git, implementasi aplikasi ArgoCD induk dan anak, konfigurasi Helm chart InvenioRDM, serta konfigurasi networking dan storage.
    
    \item \textbf{Pengujian dan Evaluasi (Mei -- Juni 2026)}
    
    Kegiatan pengujian meliputi pengujian fungsional deployment, pengujian sinkronisasi GitOps, pengujian drift detection, pengujian rollback, serta pengukuran metrik kinerja sistem.
    
    \item \textbf{Dokumentasi dan Penulisan (Maret -- Juli 2026)}
    
    Kegiatan dokumentasi dilakukan secara paralel dengan tahapan lainnya. Penulisan proposal tugas akhir dilakukan pada semester ini, sedangkan penulisan skripsi lengkap akan dilanjutkan pada semester berikutnya.
    
    \item \textbf{Seminar Proposal / Sidang (Juli 2026)}
    
    Tahap akhir adalah presentasi dan sidang proposal tugas akhir di depan dosen penguji.
\end{enumerate}


%------------------------------------------------------------
% Bibliography
%------------------------------------------------------------
\bibliography{references}

\end{document}
